% DO CASE -- os/samir/cmd/flow.c LK_DOCASE/LK_CASE/LK_OTHERWISE/LK_ENDCASE
% (first true CASE runs; guard must be Logical, #37; program lines only).
% Gate: test-interp-flow.
\begin{XbaseEntry}{DO CASE}{compound}
\XAlso{CASE}{CASE}\XAlso{OTHERWISE}{OTHERWISE}\XAlso{ENDCASE}{ENDCASE}%
Chooses one of several groups of commands in a program.
\Syntax
\SyntaxLine{\kw{DO CASE}}
\SyntaxLine{\kw{CASE} \xp{condition}}
\SyntaxLine{\ \ \xp{commands}}
\SyntaxLine{\opt{\kw{CASE} \xp{condition}}}
\SyntaxLine{\ \ \opt{\xp{commands}}}
\SyntaxLine{\rep}
\SyntaxLine{\opt{\kw{OTHERWISE}}}
\SyntaxLine{\ \ \opt{\xp{commands}}}
\SyntaxLine{\kw{ENDCASE}}
\Usage
InitechBase tests the condition of each \cmd{CASE} in turn.  It carries out
the commands of the first \cmd{CASE} whose condition is true, and then goes on
after \cmd{ENDCASE}; the other \cmd{CASE} groups are skipped, even if their
conditions are also true.  If no condition is true, the commands after
\cmd{OTHERWISE} are carried out; without \cmd{OTHERWISE}, none are.

Each condition must be a logical expression.  A condition of another type
stops the program with \xmsg{37}{Not a Logical expression.}

\cmd{DO CASE} is used only in programs.  The commands in each group may be any
commands, including further \cmd{DO CASE}, \cmd{IF} and \cmd{DO WHILE}
structures.
\Tips
Use \cmd{DO CASE} instead of a chain of \cmd{IF} commands whenever there are
more than two possibilities.
\Examples
The program \file{DEPTNAME.PRG} gives the name of the department of the
current employee:
\ProgramListing{DEPTNAME.PRG}
\DotExample{examples/c-docase}
\SeeAlso{IF, DO WHILE}
\end{XbaseEntry}

